\section{When projection is a physical correction}
Projection gives a unique representative in the selected parity class under
its metric. To interpret the removed component as an error one additionally
needs a qualified reflection model, matched magnetic state, aligned axes,
and an uncertainty budget that includes interpolation. Genuine construction
asymmetry or loss-state differences may be removed by the same operator.
Symmetry cannot distinguish those contributions from sensor error.

Independent calibration can justify a physical correction; symmetry alone
justifies a model projection and residual. Resistance-equivalent normalization
and geometric error fingerprints are conditional hypotheses in the companion
identifiability study \cite{companion-diagnostics}, not additional observations.

For an explicit sign regression, the generated synthetic check injects
$\varepsilon_R=\hat R_s-R_s=\SI{0.02}{\ohm}$ at
$\omega_e=\SI{1000}{\per\second}$. It reproduces
$r_d=-\varepsilon_R i_q/\omega_e$ and
$r_q=+\varepsilon_R i_d/\omega_e$ to $3.4\times10^{-21}$ Wb. Replacing the
resistance mismatch by the current-proportional voltage error
$\vect\varepsilon_u=-\varepsilon_R\current$ produces the identical reconstructed
flux perturbation to the same tolerance. Thus the resistance-equivalent pattern
is reproducible, including its sign, but resistance itself is not identifiable
from symmetry.
